% Jawaban bernomor ganjil yang terhubung langsung dengan defining_probability.tex:
% hanya latihan 1, 3, 5, 7, 9, dan 11. Jawaban sumber adalah materi lapisan teks
% CC BY-SA; tidak ada solusi instruktur terbatas yang ditambahkan.
\eocesolch{Probabilitas}
\begin{multicols}{2}

% 1
\eocesol{(a)~Salah. Ini adalah percobaan independen.
(b)~Salah. Ada kartu bergambar berwarna merah.
(c)~Benar. Sebuah kartu tidak mungkin sekaligus kartu bergambar dan ace.}

% 3
\eocesol{(a)~10 lemparan. Lebih sedikit lemparan berarti proporsi kepala dalam sampel lebih bervariasi,
sehingga peluang memperoleh setidaknya 60\% kepala lebih besar.
(b)~100 lemparan. Lebih banyak lemparan membuat proporsi kepala yang diamati sering lebih dekat
ke rata-rata, 0.50, dan karena itu juga di atas 0.40.
(c)~100 lemparan. Dengan lebih banyak lemparan, proporsi kepala yang diamati sering lebih dekat
ke rata-rata, 0.50.
(d)~10 lemparan. Lebih sedikit lemparan meningkatkan variasi fraksi lemparan yang menghasilkan kepala.}

% 5
\eocesol{(a)~$0.5^{10}$ = 0.00098.
(b)~$0.5^{10}$ = 0.00098.
(c)~$P$(setidaknya satu ekor) = $1 - P$(tidak ada ekor) = $1 - (0.5^{10}) \approx 1 - 0.001 = 0.999$.}

% 7
\eocesol{(a)~Tidak, ada pemilih yang sekaligus Independen dan pemilih mengambang. \\
(b)\\
\FigureFullPath[Diagram Venn untuk variabel “Independen” dan “Mengambang” ditampilkan; kedua lingkaran sebagian beririsan. Bagian lingkaran “Independen” yang tidak beririsan berlabel “24”. Bagian lingkaran “Mengambang” yang tidak beririsan berlabel “12”. Bagian irisan berlabel “11”.]{0.25}{ch_probability/figures/eoce/swing_voters/swing_voters.pdf} \\
(c)~Setiap pemilih Independen adalah pemilih mengambang atau bukan. Karena 35\% pemilih
adalah Independen dan 11\% adalah keduanya, 24\% sisanya pasti bukan pemilih mengambang.
(d)~0.47.
(e)~0.53.
(f)~P(Independen) $\times$ P(mengambang) = $0.35\times0.23 = 0.08$, yang
tidak sama dengan P(Independen dan mengambang) = 0.11, sehingga kedua peristiwa dependen.}

\end{multicols}
\begin{multicols}{2}

% 9
\eocesol{(a)~Jika kelas tidak dinilai berdasarkan kurva, keduanya independen. Jika dinilai
berdasarkan kurva, keduanya tidak independen maupun saling lepas--kecuali pengajar hanya akan
memberikan satu nilai A, situasi yang kita abaikan pada bagian~(b) dan~(c).
(b)~Keduanya mungkin tidak independen: jika Anda belajar bersama, kebiasaan belajar Anda
akan berkaitan, sehingga kinerja kuliah juga mungkin berkaitan.
(c)~Tidak. Lihat jawaban bagian~(a) ketika mata kuliah tidak dinilai berdasarkan kurva.
Secara umum, jika dua hal tidak berkaitan (independen), terjadinya salah satu tidak menghalangi
terjadinya yang lain.}

% 11
\eocesol{(a)~$0.16 + 0.09 = 0.25$.
(b)~$0.17 + 0.09 = 0.26$.
(c)~Dengan asumsi tingkat pendidikan suami dan istri independen:
$0.25 \times 0.26 = 0.065$. Perhatikan bahwa kita sebenarnya juga membuat asumsi kedua:
keputusan untuk menikah tidak berkaitan dengan tingkat pendidikan.
(d)~Asumsi independensi suami/istri mungkin tidak wajar karena orang sering menikah dengan
orang lain yang tingkat pendidikannya sebanding. Apakah asumsi kedua pada bagian~(c) wajar,
silakan Anda pertimbangkan.}

\end{multicols}
